\section{Limitations}
\label{sec:limitations}

We state the benchmark's limitations plainly; several are inherent to the design
choices and are documented here rather than papered over.

\paragraph{Two domains.} v1 covers banking and customer success only. Agent
performance varies significantly by domain, so the composite should not be read as
a general agentic-capability number. Domain expansion is the planned grant-funded
scale-up, not a v1 claim.

\paragraph{Synthetic scenarios, public on GitHub.} Scenarios are synthetic, which
is a strong native contamination defense (there is no scraped dataset to leak), but
they do not capture the full messiness of real interactions, and the public corpus
sits on GitHub where a future model could train on it. The private holdout caps
overfitting \emph{detection}, not its possibility; the public-versus-holdout gap is
a tripwire, not a guarantee.

\paragraph{Author-model concentration.} The generation pipeline supports
multi-author generation across labs and enforces the author-is-never-a-contestant
guard, but the committed public corpus is heavily skewed to one author: 76 of the
92 scenarios (83\%) carry \texttt{author\_model} \texttt{claude-opus-4-8}, with
the remainder authored by Kimi K2.6 (7), GLM-4.6 (5), and human writers (4). The
dominant author model is not under test, so the contamination guard is satisfied,
but it is a sibling snapshot of the Claude Opus panel judge and shares a lab with
the two Claude contestants. Any resulting familiarity effect on judge or
contestant behavior is unmeasured; the published per-judge deltas and the
\texttt{same\_lab\_check} measure judge-side favoritism but not author-side style
bias. Diversifying authorship is the priority for corpus growth.

\paragraph{Effectively cooperative user simulation.} Scenario personas vary in
tone and include adversarial personas, but v1 uses a single simulator model with
prompt-only persona text and no behaviorally calibrated profile set. The published
evidence is that off-the-shelf LLM user simulators skew over-cooperative and
stylistically uniform relative to real users regardless of persona
instructions~\cite{lostinsim2026,sim2real2026,realusersim2026}, which inflates
measured agent performance and hides failure modes that only impatient, confused,
or adversarial users trigger. Behavioral simulator profiles (impatient,
technically-confused, adversarial) with persona-stratified reporting are
implemented in the harness; the published leaderboard still scores every model
under the cooperative default, with profile-stratified robustness tables
reported separately, and the adversarial-simulator judge audit remains future
work (see below).

\paragraph{Single-session, ten-turn horizon.} Every scenario is one conversation
capped at ten user turns. The benchmark therefore measures within-conversation
reliability and recovery, not long-horizon coherence over extended autonomous
operation~\cite{vendingbench2025}, and not memory persistence across sessions.
For deployed banking and customer-success agents, whose sessions are short, this
is arguably the more deployment-relevant regime, but the composite should not be
read as evidence about long-horizon autonomy.

\paragraph{LLM-as-judge for the soft dimensions.} Two thirds of efficacy is judge
scored, and judge quality is bounded by the judge models' capabilities. The
multi-judge panel, median consensus, per-scenario atomic rubric criteria,
published per-judge deltas, length-bias
regression, and calibration protocol mitigate this, but they do not remove the
underlying dependence. State verification, the judge-independent third, is the
direct hedge.

\paragraph{Tool simulation generates the response surface.} For v0.2 scenarios the
underlying \emph{values} are pinned to a canonical ground-truth world the simulator
must answer from and mutate, and the resulting state is verified deterministically,
so balances, IDs, and records stay coherent. But the simulator still generates the
response \emph{surface} with an LLM, so edge cases in real API behavior may not be
perfectly represented.

\paragraph{Unseeded simulators.} The user and tool simulators run unseeded and the
user simulator at temperature above zero, so runs are not bit-for-bit reproducible;
the pre-registration says so explicitly. The agent under test runs at temperature
0.0, which aids reproducibility but misses failure modes that appear only at higher
temperatures. A single LLM plays the user simulator, and which model plays the user
can swing agent success by up to roughly nine points~\cite{lostinsim2026}; the
completion-decoupling design removes the simulator from the success-declaration
path, and a planned sim-sensitivity check will quantify the residual swing, but the
single-simulator risk is real until that check runs.

\paragraph{Field-relative composite.} CLEAR is min-max normalized across the
evaluated models, so adding or removing a model rescales every other model's
composite. Cross-run comparison must use raw per-dimension scores, and a single-run
ranking is only as stable as its roster.

\paragraph{Cost approximation.} Token costs use published per-token pricing; actual
costs vary with volume discounts and cached tokens, and OpenRouter slugs pin the
model but not the serving provider or quantization (recorded per run, but variable).

\paragraph{Single maintainer.} The benchmark is run by one person plus community
contributions. The governance policy (no silent retraction, pre-registration,
pinning, automated weekly runs) is designed precisely to make a single-maintainer
operation auditable, but it does not substitute for the institutional backing the
larger benchmarks have.

\subsection*{Future work}

Several of the limitations above have committed mitigations, tracked as open
issues in the repository (\#57--\#60) so the roadmap is itself auditable. Each
would be a version bump under the policy of Section~\ref{sec:versioning}.
\emph{A recovery-probe scenario tier} (\#57) injects a deliberate error or
ambiguity mid-conversation and verifies recovery deterministically, bringing the
goal-maintenance question inside the ten-turn
horizon~\cite{vendingbench2025}. \emph{Dual-control scenarios} (\#58) let the
simulated user also act on the shared world, the
$\tau^2$-bench pattern~\cite{tau2bench2025}, which is the dominant real-world
complexity the single-control design misses. \emph{An adversarial-simulator
audit pass} (\#59) runs a subset under the implemented adversarial profile and
uses rubric-versus-adversarial divergence as a judge-refinement
signal~\cite{realusersim2026}. \emph{Parameterized scenario templates} (\#60)
resample surface entities (names, amounts, dates) per run so a model that has
seen the public corpus cannot lean on memorized surface features. None of these
are v1 claims; they are listed so the gap between what is measured now and what
should be measured next is explicit.
